\documentclass[12pt,letterpaper]{article}
\usepackage[margin=1in]{geometry}
\usepackage[T1]{fontenc}
\usepackage[utf8]{inputenc}
\usepackage[french]{babel}
\usepackage{graphicx} % Required for inserting images
\usepackage{url}
\usepackage{etoolbox}
\usepackage{xcolor}
\usepackage{amsmath}
\usepackage{listings}

% Code listing style for Python
\lstdefinestyle{pythonstyle}{
  language=Python,
  basicstyle=\ttfamily\small,
  keywordstyle=\color{blue!70!black}\bfseries,
  commentstyle=\color{green!40!black}\itshape,
  stringstyle=\color{orange!70!black},
  numbers=left,
  numberstyle=\scriptsize\color{gray!70},
  stepnumber=1,
  numbersep=10pt,
  frame=single,
  framerule=0.3pt,
  rulecolor=\color{gray!40},
  backgroundcolor=\color{gray!5},
  breaklines=true,
  breakatwhitespace=true,
  showstringspaces=false,
  tabsize=4,
  columns=fullflexible,
  keepspaces=true,
  xleftmargin=12pt,
  xrightmargin=6pt,
  aboveskip=8pt,
  belowskip=8pt,
  captionpos=b
}

\lstset{style=pythonstyle}


% Start each section on a new page
\pretocmd{\section}{\clearpage}{}{}

% Remove section numbering
\setcounter{secnumdepth}{0}


\title{MTH8302~- Devoir 1: Rappel Mathématique}
\author{Leonard Excoffier}
\date{Hiver 2026}

\begin{document}

\maketitle

\tableofcontents

\section*{Introduction}
Le dépôt GitHub~\cite{repo} contient le code complet, structuré et prêt à être exécuté; il est plus simple de l'utiliser directement plutôt que de copier-coller les extraits.

\section{Problème 1: Gradients et Hessiennes}
\subsection*{Exemple de code Python}
\lstinputlisting[caption={Exemple Python (main.py)}]{../src/main.py}

\subsection{Question 1}
\subsection{Question 2}
\subsection{Question 3}
\subsection{Question 4}
\subsection{Question 5}
\subsection{Question 6}

\section{Problème 2: Optimisation, Descentes de Gradient vs Newton et analyse spectrale}
\subsection{a}
\subsection{b}
\subsection{c}
\subsection{d}
\subsection{e}
\subsection{f}
\subsection{g}

\section{Problème 3: Estimation par la Méthode du Maximum de Vraisemblance et Probabilités des Estimateurs}
\subsection{a}
\subsection{b}
\subsection{c}
\subsection{d}

\section{Problème 4: Test F de Fisher et Intervalle de Confiance}

\subsection{a}
\textbf{Hypothèses.} On teste
\(H_0: \sigma_A^2 = \sigma_B^2\) contre \(H_1: \sigma_A^2 \neq \sigma_B^2\) (test bilatéral), avec \(\alpha = 0.05\).

Les variances d'échantillon calculées sont \(s_A^2 = 16.988889\) et \(s_B^2 = 10.266667\), d'où la statistique \(F_{\text{stat}} = s_A^2 / s_B^2 = 1.654762\).
Les valeurs critiques pour \((\nu_1, \nu_2) = (9, 9)\) sont \(F_{0.025} = 0.248386\) et \(F_{0.975} = 4.025994\), et la p-valeur vaut \(0.464710\).
Comme \(F_{\text{stat}}\) est dans la région d'acceptation et \(p > 0.05\), on ne rejette pas \(H_0\): il n'y a pas d'évidence d'une différence de variances au seuil de 5\%.

\subsection{b}
  \begin{figure}[h]
  \centering
  \includegraphics[width=0.85\textwidth]{../output/4.png}
  \caption{Distribution \(F(9,9)\) avec les zones critiques et la statistique observée \(F_{\text{stat}}\).}%
  \label{fig:fisher-f}
\end{figure}

La figure~\ref{fig:fisher-f} montre la densité \(F(9,9)\), les régions critiques (en rouge) et la position de \(F_{\text{stat}}\), qui se situe hors des zones de rejet.
\subsection{c}
\textbf{Intervalle de confiance à 95\% pour \(\sigma_A^2 / \sigma_B^2\).}
En utilisant les quantiles \(F_{0.025}\) et \(F_{0.975}\), l'intervalle est
\[
\left[\frac{F_{\text{stat}}}{F_{0.975}}, \frac{F_{\text{stat}}}{F_{0.025}}\right]=[0.411019,\; 6.662062].
\]

\subsection{d}
\textbf{Code et sortie.} Le code source est fourni au listing~\ref{lst:fisher-code} et la sortie correspondante au listing~\ref{lst:fisher-output}.

\lstinputlisting[caption={Code Python pour le test F et l'intervalle de confiance (src/4.py)}, label={lst:fisher-code}]{../src/4.py}

\lstinputlisting[caption={Sortie du programme (output/4.txt)}, label={lst:fisher-output}, numbers=none, basicstyle=\ttfamily\small]{../output/4.txt}

\textbf{Interprétation.} L'intervalle de confiance contient la valeur \(1\), donc un ratio \(\sigma_A^2 / \sigma_B^2 = 1\) est plausible.
Cette conclusion est cohérente avec le test F qui ne rejette pas \(H_0\).


\section{Problème 5: Loi des Grands Nombres et Théorème Central Limite}
\subsection{Question 1}
\subsection{Question 2}
\subsection{Question 3}

\begin{thebibliography}{9}
\bibitem{repo}
Leonard Excoffier.~\emph{MTH8302~- Devoir 1: Rappel Mathématique}.\ GitHub repository.
\url{https://github.com/excoffierleonard/mth8302-devoir1.git}. Consulté le 18 février 2026.
\end{thebibliography}

\end{document}
